% Figura redibujada a partir de "Imagenes/Estructuras.png" (Tema 1).
% Coordenadas en píxeles del PNG original (616 x 451 px), escaladas para que
% ocupe lo mismo que el PNG en el documento: 193.5 pt x 141.7 pt (\linewidth
% de la wrapfigure). Escala de energías: y = 8 + 60.83 (4 - E).
\documentclass[12pt,border=0pt]{standalone}
\usepackage[spanish,es-tabla]{babel}
\usepackage[utf8]{inputenc}
\usepackage{amsmath}
\usepackage{xcolor}
\usepackage{tikz}
\usepackage{pgfplots}
\usepackage{circuitikz}
\usetikzlibrary{babel, circuits.ee}
\usetikzlibrary{calc}

\definecolor{applegreen}{rgb}{0.55, 0.71, 0.0}
\definecolor{airforceblue}{rgb}{0.36, 0.54, 0.66}
\definecolor{amber}{rgb}{1.0, 0.49, 0.0}

\newcommand{\diff}[3][1]{%
	\IfBlankTF{#1}{\frac{\mathrm{d}{#2}}{\mathrm{d}{#3}}}%
	{\frac{\mathrm{d}^{#1} {#2}}{\mathrm{d}{#3}^{#1}}}%
}%
\newcommand{\pdiff}[3][1]{%
	\IfBlankTF{#1}{\frac{\partial {#2}}{\partial {#3}}}%
	{\frac{\partial^{#1} {#2}}{\partial {#3}^{#1}}}%
}%

% Marco de un panel: #1 desplazamiento horizontal, #2 letra del panel.
\newcommand{\panel}[2]{%
    \begin{scope}[shift={(#1,0)}]
        \draw (37,8) rectangle (287,373);
        \draw (160,8) -- (160,373);
        \foreach \e in {-2,...,4} {
            \draw (37,{yE(\e)}) -- (44,{yE(\e)});
            \node[anchor=east, inner sep=1.5pt] at (37,{yE(\e)}) {$\e$};
        }
        \node[rotate=90] at (10,190) {Energía (eV)};
        \node at (105,386) {[111]};
        \node at (160,386) {$0$};
        \node at (225,386) {[100]};
        \node at (165,408) {$k$};
        \node at (160,436) {(#2)};
    \end{scope}}
% Banda de valencia (igual en los dos paneles, máximo en k = 0).
\newcommand{\bandavalencia}{%
    \draw[ultra thick] plot[smooth, tension=0.55] coordinates
        {(37,{yE(-1.27)}) (70,{yE(-1.05)}) (100,{yE(-0.72)}) (130,{yE(-0.3)}) (160,{yE(0)})
         (178,{yE(-0.2)}) (193,{yE(-0.7)}) (207,{yE(-1.3)}) (222,{yE(-2)})};}
% Electrones (puntos rellenos) y huecos (círculos vacíos) centrados en #1.
\newcommand{\electrones}[1]{%
    \foreach \dx/\dy in {-5/-4, 4/-5, -4/5, 5/4}
        \fill[blue] ($#1 + (\dx,\dy)$) circle (4.2);}
\newcommand{\huecos}[1]{%
    \foreach \dx/\dy in {-5/-4, 4/-5, -4/5, 5/4}
        \draw[blue, fill=white, line width=0.3pt] ($#1 + (\dx,\dy)$) circle (3.8);}

\begin{document}
\begin{tikzpicture}[x=0.314089pt, y=-0.314089pt, font=\tiny,
    declare function={yE(\e) = 8 + 60.833*(4 - \e);}]
    \useasboundingbox (0,0) rectangle (616,451);
    % ================= (a) GaAs =================
    \panel{0}{a}
    % Banda de conducción: valle L (izquierda), Gamma (centro), X (derecha)
    \draw[ultra thick] plot[smooth, tension=0.55] coordinates
        {(37,{yE(1.75)}) (55,{yE(1.78)}) (85,{yE(2.12)}) (122,{yE(2.5)}) (140,{yE(2.2)})
         (152,{yE(1.6)}) (160,{yE(1.43)}) (170,{yE(1.65)}) (183,{yE(2.6)}) (200,{yE(3.35)})
         (215,{yE(2.95)}) (240,{yE(2.2)}) (262,{yE(2.0)}) (287,{yE(1.95)})};
    \bandavalencia
    % Delta E entre los valles L y Gamma
    \draw (37,{yE(1.75)}) -- (120,{yE(1.75)});
    \draw (75,{yE(1.43)}) -- (160,{yE(1.43)});
    \draw[-latex] (100,{yE(2.35)}) -- (100,{yE(1.75)});
    \draw[-latex] (100,{yE(0.83)}) -- (100,{yE(1.43)});
    \node[fill=white, inner sep=0.3pt, font=\fontsize{5}{6}\selectfont] at (89,{yE(1.59)}) {$\Delta E = 0.31$};
    % Gap directo
    \draw[latex-latex, blue] (138,{yE(1.41)}) -- (138,{yE(0.03)});
    \node[blue] at (121,{yE(0.72)}) {$E_\mathrm{g}$};
    \electrones{(158,{yE(1.68)})}
    \huecos{(158,{yE(-0.22)})}
    \node at (82,22) {GaAs};
    \node[align=center] at (226,32) {Banda de\\conducción};
    \node[align=center] at (112,352) {Banda de\\valencia};
    % ================= (b) Si =================
    \panel{325}{b}
    \begin{scope}[shift={(325,0)}]
        % Banda de conducción con el mínimo en la dirección [100]
        \draw[ultra thick] plot[smooth, tension=0.55] coordinates
            {(37,{yE(2.15)}) (62,{yE(2.2)}) (92,{yE(2.45)}) (122,{yE(2.92)}) (145,{yE(3.25)})
             (160,{yE(3.36)}) (185,{yE(3.1)}) (205,{yE(2.6)}) (228,{yE(1.85)}) (252,{yE(1.25)})
             (270,{yE(1.07)}) (287,{yE(1.1)})};
        \bandavalencia
        % Gap indirecto
        \draw (227,{yE(1.07)}) -- (282,{yE(1.07)});
        \draw[dashed, blue] (268,{yE(1.0)}) -- (268,373);
        \draw[latex-latex, blue] (241,{yE(1.05)}) -- (241,{yE(0.03)});
        \node[blue] at (229,{yE(0.52)}) {$E_\mathrm{g}$};
        \electrones{(266,{yE(1.24)})}
        \huecos{(160,{yE(-0.22)})}
        \node at (82,22) {Si};
        \node[align=center] at (226,32) {Banda de\\conducción};
        \node[align=center] at (112,352) {Banda de\\valencia};
    \end{scope}
\end{tikzpicture}
\end{document}
